\clearpage\thispagestyle{empty}
\pdfbookmark[0]{Quellen zum Pallas-Bericht vom August 1802}{pallas180208-sources}
\begin{center}\textbf{Quellen zum Pallas-Bericht vom August 1802}\end{center}
\noindent Diese KI-Übertragung in heutiges Deutsch folgt dem deutschen Auszug in Carl Friedrich Gauß’ \href{https://gdz.sub.uni-goettingen.de/id/PPN236007467}{\textit{Werken}, Band~VI} (Göttingen, 1874), S.~221--222. Franz Xaver von Zach berichtet über Pallas und zitiert Carl Friedrich Gauß und Heinrich Wilhelm Matthias Olbers.
\par\smallskip\noindent Von Zachs \href{https://api.digitale-sammlungen.de/iiif/image/v2/bsb10538598_00195/full/1450,/0/default.jpg}{ursprünglicher Aufsatz}, Nr.~XVIII, erschien in der \textit{Monatlichen Correspondenz zur Beförderung der Erd- und Himmels-Kunde} (Gotha, Becker), Band~VI, August 1802, S.~187--196. Der Auszug beginnt auf \href{https://api.digitale-sammlungen.de/iiif/image/v2/bsb10538598_00200/full/1450,/0/default.jpg}{S.~192} mit den Meldungen aus Vilna, Prag, Kremsmünster und Padua. Der \href{https://api.digitale-sammlungen.de/iiif/image/v2/bsb10538598_00201/full/1450,/0/default.jpg}{Mailänder Vergleich} verweist auf die \href{https://api.digitale-sammlungen.de/iiif/image/v2/bsb10538598_00079/full/1450,/0/default.jpg}{Beobachtungen im Juliheft, S.~73}. Die Fortsetzung der \href{https://api.digitale-sammlungen.de/iiif/image/v2/bsb10538598_00091/full/1450,/0/default.jpg}{Ephemeride} aus dem Juliheft steht im Zeitschriftendruck auf \href{https://api.digitale-sammlungen.de/iiif/image/v2/bsb10538598_00204/full/1450,/0/default.jpg}{S.~196}. An ihrer Stelle steht in den \textit{Werken} eine Bemerkung des Herausgebers in eckigen Klammern.
\par\smallskip\noindent Die Zeitschriftentabellen enthalten keine Spalten für die mittlere Zeit. Für den letzten Deklinationsunterschied im Mailänder Vergleich steht in der Zeitschrift \(+1{,}4\), in den \textit{Werken} \(-1.4\). In der \href{https://api.digitale-sammlungen.de/iiif/image/v2/bsb10538598_00202/full/1450,/0/default.jpg}{Bremer Tabelle} ist die zweite berechnete Spalte in der Zeitschrift mit nördlicher Deklination überschrieben. In den \textit{Werken} sind beide berechneten Spalten mit Rektaszension überschrieben. Die Übertragungen folgen jeweils den \textit{Werken}.
\par\smallskip\noindent \href{https://doi.org/10.1098/rstl.1802.0010}{William Herschels Aufsatz}, S.~216--217, beschreibt seine Beobachtungen vom 22. April. Eine eigenständige englische Übertragung eines Teils der Erörterung der Durchmesser findet sich bei \href{https://doi.org/10.1007/978-3-319-07614-0_14}{Clifford J.\ Cunningham und Wayne Orchiston, Abschnitt~6.2}.
\par\smallskip\noindent Übertragung und Satz wurden mit KI erstellt. Die Digitalisate der Zeitschrift stammen von der Bayerischen Staatsbibliothek.
\clearpage
